\chapter{Resultados e Discussão} \label{chap:resultados}

% Objetivo: apresentar resultados, comparações e interpretações com base em
% evidências rastreáveis. Status: aguarda a rodada confirmatória; as seções
% abaixo são espaços reservados com o conteúdo que cada uma reportará.

\section{Cenários experimentais}
Os experimentos deste trabalho têm duas naturezas distintas. A \textbf{rodada exploratória} consiste na busca de hiperparâmetros do TFT, avaliada exclusivamente pela perda no bloco \texttt{early\_stop} de cada \textit{fold}, sem persistência de previsões fora da amostra; ela não constitui evidência e serve apenas para fixar a configuração do candidato. A \textbf{rodada confirmatória} é executada uma única vez, com hiperparâmetros congelados, sobre o cohort pré-registrado (sementes \(\times\) \textit{folds}), contra a hierarquia completa de \textit{baselines}, sob alinhamento estrito por \textit{target timestamp}. As seções seguintes são organizadas como espaços reservados que serão preenchidos com os artefatos da rodada confirmatória.

\section{Calibração por horizonte (H1)}
\textit{A ser preenchido após a rodada confirmatória.} Esta seção reportará, por horizonte (\(h+1\), \(h+7\); \(h+30\) suplementar), a cobertura marginal em cada nível da grade e o \textit{reliability diagram}, o PICP e o MPIW dos intervalos definidos no pré-registro, o \textit{interval score} e os testes de cobertura condicional de Christoffersen e Kupiec, confrontados com as bandas pré-registradas e com os intervalos conformais (CQR) calibrados no bloco \texttt{calib}. A taxa de degeneração da grade e a proporção de previsões alteradas pelo rearranjo serão reportadas por execução.

\section{\textit{Skill} relativo contra a hierarquia de \textit{baselines} (H2)}
\textit{A ser preenchido após a rodada confirmatória.} Para os horizontes aprovados em H1, esta seção reportará a \textit{pinball loss} média e o CRPS do candidato e de cada \textit{baseline} (retorno nulo, média histórica, AR(1), EWMA-vol, quantis empíricos históricos e GBM quantílico), com a dispersão entre sementes e \textit{folds}.

\section{Análise estatística pareada}
\textit{A ser preenchido após a rodada confirmatória.} A análise pareada utilizará o teste de Diebold-Mariano \cite{DIEBOLDMARIANO1995} unilateral sobre a \textit{pinball loss}, com variância HAC de defasagem \(h-1\) e correção HLN \cite{HARVEY1997}, ajuste de Holm \cite{HOLM1979} sobre a família de comparações e o \textit{Model Confidence Set} \cite{HANSEN2011MCS}, que integra o \textit{scorecard} como leitura complementar. O veredito de H1 e H2 por horizonte será o produzido mecanicamente pelo \textit{scorecard} (Seção~\ref{sec:scorecard}).

\section{Perfil descritivo pontual}
\textit{A ser preenchido após a rodada confirmatória.} RMSE, MAE e acurácia direcional por horizonte, reportados como perfil descritivo, sem papel no veredito.

\section{Contribuição de famílias (H3)}
\textit{A ser preenchido após a rodada confirmatória.} A contribuição relativa das famílias (preço, técnico, sentimento, fundamento) será reportada por horizonte segundo os três métodos --- pesos da VSN, importância por permutação com intervalo \textit{bootstrap} e ablação ---, destacando as famílias com consistência em pelo menos dois métodos.

\section{Ameaças à validade}
As principais ameaças à validade são: (i) sensibilidade dos resultados ao recorte de ativo, horizonte e janela temporal; (ii) possível subpotência estatística para detectar diferenças pequenas entre o candidato e os comparadores fortes; (iii) interpretação indevida quando comparações são feitas fora da mesma interseção temporal; e (iv) não permutabilidade das séries temporais, que limita a garantia teórica da predição conformal --- razão pela qual sua cobertura é reportada como empírica \cite{BARBER2023}. Para mitigar esses riscos, a análise adota coorte congelada, alinhamento temporal estrito, correção de múltiplas comparações e leitura conjunta de métricas de calibração, \textit{skill} e testes pareados.

A generalização externa dos achados depende de replicação em outros ativos e períodos. As conclusões devem ser entendidas como evidência para o escopo experimental definido, e não como afirmação universal sobre o desempenho relativo de arquiteturas ou famílias de \textit{features}.
